\documentclass[UTF8]{ctexart}
\usepackage{amsmath,amssymb}
\usepackage[a4paper,margin=2.5cm]{geometry}
\xeCJKDeclareCharClass{CJK}{"2460 -> "2473} % 让 ①②③ 等圈号使用中文字体
\title{几何与代数 0921}
\date{}

\begin{document}
\maketitle

\subsection*{0.\ Recall：向量的外积}

$\vec a\times\vec b$ 定义为长度为
$|\vec a||\vec b|\sin\langle\vec a,\vec b\rangle$、方向满足
①与 $\vec a,\vec b$ 垂直，②$(\vec a,\vec b,\vec a\times\vec b)$
是右手系的向量。

\subsubsection*{几何意义}

\begin{itemize}
\item $|\vec a\times\vec b|$ 是平行四边形的面积；
\item $\vec e\times\vec a=\vec e\times\vec a_h$，
即右旋 $\dfrac\pi2$。
\end{itemize}

\subsubsection*{外积性质}

\begin{itemize}
\item 反交换律\quad
$\vec a\times\vec b=-\vec b\times\vec a$；
\item 结合律\quad
$(\lambda\vec a)\times\vec b=\lambda(\vec a\times\vec b)$；
\item 分配律\quad
$\vec e\times(\vec a+\vec b)
=\vec e\times\vec a+\vec e\times\vec b$，
左乘同理。
\end{itemize}

\subsubsection*{坐标表示}

在右手直角坐标系
$[O;\vec e_1,\vec e_2,\vec e_3]$ 中，
$\vec a=(x_1,y_1,z_1)$，$\vec b=(x_2,y_2,z_2)$，则
\[
\vec a\times\vec b
=(y_1z_2-y_2z_1,\ z_1x_2-z_2x_1,\ x_1y_2-x_2y_1).
\]

\subsection*{1.\ 行列式}

\subsubsection*{二阶行列式}

\[
\begin{cases}
a_{11}x+a_{12}y=b_1,\\
a_{21}x+a_{22}y=b_2.
\end{cases}
\]
假设方程有唯一解
$\Longleftrightarrow a_{11}a_{22}\neq a_{12}a_{21}$，
有如下的解：
\[
x=\frac{b_1a_{22}-b_2a_{12}}{a_{11}a_{22}-a_{12}a_{21}},\qquad
y=\frac{a_{11}b_2-a_{21}b_1}{a_{11}a_{22}-a_{12}a_{21}} .
\]
记作二阶行列式如下
\[
\begin{vmatrix}
a_{11} & a_{12}\\
a_{21} & a_{22}
\end{vmatrix}
\stackrel{\Delta}{=}a_{11}a_{22}-a_{21}a_{12},
\]
上面的解就可以记作
\[
x=\frac{
\begin{vmatrix}
b_1 & a_{12}\\
b_2 & a_{22}
\end{vmatrix}
}{
\begin{vmatrix}
a_{11} & a_{12}\\
a_{21} & a_{22}
\end{vmatrix}
},\qquad
y=\frac{
\begin{vmatrix}
a_{11} & b_1\\
a_{21} & b_2
\end{vmatrix}
}{
\begin{vmatrix}
a_{11} & a_{12}\\
a_{21} & a_{22}
\end{vmatrix}
}.
\]

\subsubsection*{三阶行列式}

扩展到三阶行列式，有三元一次方程
\[
\begin{cases}
a_{11}x+a_{12}y+a_{13}z=b_1,\\
a_{21}x+a_{22}y+a_{23}z=b_2,\\
a_{31}x+a_{32}y+a_{33}z=b_3 .
\end{cases}
\]
下定义三阶行列式
\[
\begin{vmatrix}
a_{11} & a_{12} & a_{13}\\
a_{21} & a_{22} & a_{23}\\
a_{31} & a_{32} & a_{33}
\end{vmatrix}
\stackrel{\Delta}{=}
a_{11}\begin{vmatrix}
a_{22} & a_{23}\\
a_{32} & a_{33}
\end{vmatrix}
-a_{12}\begin{vmatrix}
a_{21} & a_{23}\\
a_{31} & a_{33}
\end{vmatrix}
+a_{13}\begin{vmatrix}
a_{21} & a_{22}\\
a_{31} & a_{32}
\end{vmatrix}.
\]
有了三阶行列式，我们就可以写方程组的解了
\[
x=\frac{
\begin{vmatrix}
b_1 & a_{12} & a_{13}\\
b_2 & a_{22} & a_{23}\\
b_3 & a_{32} & a_{33}
\end{vmatrix}
}{
\begin{vmatrix}
a_{11} & a_{12} & a_{13}\\
a_{21} & a_{22} & a_{23}\\
a_{31} & a_{32} & a_{33}
\end{vmatrix}
},\qquad
y,z\ \text{类似}.
\]

借用行列式，我们可以把叉乘写成如下模式
\[
\vec a\times\vec b=
\begin{vmatrix}
x_1 & x_2 & \vec e_1\\
y_1 & y_2 & \vec e_2\\
z_1 & z_2 & \vec e_3
\end{vmatrix}
\]
也可以写成线性组合的形式
\[
=\begin{vmatrix}
y_1 & y_2\\
z_1 & z_2
\end{vmatrix}\vec e_1
-\begin{vmatrix}
x_1 & x_2\\
z_1 & z_2
\end{vmatrix}\vec e_2
+\begin{vmatrix}
x_1 & x_2\\
y_1 & y_2
\end{vmatrix}\vec e_3 .
\]

\subsubsection*{叉乘理解直线和坐标面平行}

直线 $L$
\[
\begin{cases}
A_1x+B_1y+C_1z+D_1=0,\\
A_2x+B_2y+C_2z+D_2=0,
\end{cases}
\]
则 $L$ 的方向向量为
\[
\vec v=\vec n_1\times\vec n_2
=\begin{vmatrix}
A_1 & A_2 & \vec e_1\\
B_1 & B_2 & \vec e_2\\
C_1 & C_2 & \vec e_3
\end{vmatrix}
=(B_1C_2-B_2C_1,\ -A_1C_2+A_2C_1,\ A_1B_2-A_2B_1).
\]
与 $xy$ 平面平行即 $A_1B_2-A_2B_1=0$。

\subsubsection*{点到直线的距离}

有直线
\[
L:\ \frac{x-x_0}{l}=\frac{y-y_0}{m}=\frac{z-z_0}{n},
\]
又有
\[
d=|(\overrightarrow{P_0P_1})_h|,
\]
而
\[
\overrightarrow{P_0P_1}
=(\overrightarrow{P_0P_1})_{\vec e}
+(\overrightarrow{P_0P_1})_h,
\]
可知
\[
|\overrightarrow{P_0P_1}\times\vec v|
=|(\overrightarrow{P_0P_1})_h\times\vec v|
=|(\overrightarrow{P_0P_1})_h||\vec v|,
\]
\[
d=|(\overrightarrow{P_0P_1})_h|
=\frac{|\overrightarrow{P_0P_1}\times\vec v|}{|\vec v|}.
\]

\subsubsection*{点乘、叉乘与水平分量、垂直分量}

\[
|\vec a_h|=|\vec a\times\vec e|,\qquad
|\vec a_{\vec e}|=|\vec a\cdot\vec e| .
\]

\subsection*{2.\ 二重外积公式}

\[
\vec a\times(\vec b\times\vec c)
=(\vec a\cdot\vec c)\vec b-(\vec a\cdot\vec b)\vec c .
\]

\end{document}
